\section{Microcanonical enumeration of states on a finite screen}

Consider the finite realization
\[
 \mathcal H_{\rm occ}=\bigoplus_{k=0}^{3}\CC^{g_k},
 \qquad
 H_{\rm occ}|_{\CC^{g_k}}=E_k^{\rm occ} I,
\]
with levels and multiplicities
\[
 E^{\rm occ}=(0,3,6,9),\qquad g=(1,6,12,8),
 \qquad(N,E)=(12,66).
\]
The sum \(1+6+12+8=27\) makes \(\mathcal H_{\rm occ}\) a screen of
twenty-seven modes. This spectrum is the datum of a declared finite
representation; its cardinality does not suffice to identify it with any other
occurrence of \(27\) in HMT\@. The morphism relating a specific
boundary screen to \((\mathcal H_{\rm occ},H_{\rm occ})\) must specify both
the fibers and the spectral operator.
The fermionic and bosonic generating functions are
\[
 F(y,x)=\prod_k(1+yx^{E_k^{\rm occ}})^{g_k},
 \qquad
 B(y,x)=\prod_k(1-yx^{E_k^{\rm occ}})^{-g_k}.
\]

\begin{proposition}[Microcanonical cardinalities]
\label{prop:cardinalidades-microcanonicas}
\[
 [y^{12}x^{66}]F=2\,104\,494,
 \qquad
 [y^{12}x^{66}]B=259\,436\,430.
\]
\end{proposition}
\begin{proof}
In the fermionic product, choosing \(r\) of the \(g_k\) modes contributes
\(\binom{g_k}{r}y^rx^{rE_k^{\rm occ}}\). In the bosonic product, distributing \(r\)
indistinguishable occupations among \(g_k\) modes contributes
\(\binom{g_k+r-1}{r}y^rx^{rE_k^{\rm occ}}\). The finite convolution over the
four levels produces the stated integers.
\end{proof}

\section{Bose--Einstein and Fermi--Dirac distributions}

The equilibrium formulation starts from the Hamiltonian and the total
occupation number,
\begin{equation}
 H=\sum_kE_k^{\rm occ} n_k,
 \qquad
 N=\sum_k n_k,
 \qquad
 K_\mu:=H-\mu N.
 \label{eq:hamiltoniano-gran-canonico-rev3}
\end{equation}
For \(\beta>0\), the grand canonical state is
\begin{equation}
 \rho_{\beta,\mu}
 =\frac{e^{-\beta K_\mu}}{\mathcal Z_{\beta,\mu}},
 \qquad
 \mathcal Z_{\beta,\mu}
 =\operatorname{Tr}(e^{-\beta K_\mu}).
 \label{eq:estado-gran-canonico-rev3}
\end{equation}

\begin{theorem}[Grand canonical factorization and occupation laws]
\label{thm:factorizacion-gran-canonica-rev3}
For energy levels \(E_k^{\rm occ}\) with degeneracies \(g_k\), the
fermionic trace factorizes as
\begin{equation}
 \mathcal Z_{\rm FD}
 =\prod_k\bigl(1+e^{-\beta(E_k^{\rm occ}-\mu)}\bigr)^{g_k},
 \qquad
 \langle n_k\rangle_{\rm FD}
 =\frac{g_k}{e^{\beta(E_k^{\rm occ}-\mu)}+1}.
 \label{eq:factorizacion-fd-rev3}
\end{equation}
In the bosonic sector, under the condition
\(\mu<\inf_kE_k^{\rm occ}\),
\begin{equation}
 \mathcal Z_{\rm BE}
 =\prod_k\bigl(1-e^{-\beta(E_k^{\rm occ}-\mu)}\bigr)^{-g_k},
 \qquad
 \langle n_k\rangle_{\rm BE}
 =\frac{g_k}{e^{\beta(E_k^{\rm occ}-\mu)}-1}.
 \label{eq:factorizacion-be-rev3}
\end{equation}
\end{theorem}

\begin{proof}
The operator \(K_\mu\) is a sum of commuting occupation operators, so
its trace factorizes over modes. In each fermionic mode, the
occupations \(0\) and \(1\) are summed, and the corresponding factor is
\(1+u_k\), where
\(u_k=e^{-\beta(E_k^{\rm occ}-\mu)}\). In each bosonic mode, all
nonnegative occupations are summed, giving the geometric series
\((1-u_k)^{-1}\); the condition on \(\mu\) ensures \(u_k<1\). The
powers \(g_k\) incorporate the degeneracy. Finally,
\[
 \langle n_k\rangle
 =-\frac1\beta\frac{\partial}{\partial E_k^{\rm occ}}
   \log\mathcal Z
\]
yields the two occupation laws
\cite{Bose1924,Einstein1924,Fermi1926,Dirac1926,Pathria2011}.
\end{proof}

The microcanonical derivation recovers the same result in the
thermodynamic regime. For mean fermionic occupations
\(p_k=m_k/g_k\), the combinatorial entropy per level is
\begin{equation}
 S_k\simeq
 g_k\bigl[-p_k\log p_k-(1-p_k)\log(1-p_k)\bigr].
 \label{eq:entropia-fd-nivel-rev3}
\end{equation}
The variation of \(\sum_kS_k\), subject to
\(\sum_kg_kp_k=N\) and
\(\sum_kg_kE_k^{\rm occ}p_k=E\), gives
\begin{equation}
 \log\frac{1-p_k}{p_k}=\alpha+\beta E_k^{\rm occ},
 \qquad
 \mu=-\frac{\alpha}{\beta},
 \label{eq:variacion-fd-rev3}
\end{equation}
and hence the Fermi--Dirac occupation of
\eqref{eq:factorizacion-fd-rev3}. The multipliers \(\beta\) and \(\mu\)
are the dual coordinates of the energy and population constraints.

For mean bosonic occupations
\(\overline n_k=m_k/g_k\), the log-multiplicity of the level is
\begin{equation}
 S_k\simeq
 g_k\bigl[(1+\overline n_k)\log(1+\overline n_k)
          -\overline n_k\log\overline n_k\bigr].
 \label{eq:entropia-be-nivel-rev3}
\end{equation}
The variation of \(\sum_kS_k\), subject to
\(\sum_kg_k\overline n_k=N\) and
\(\sum_kg_kE_k^{\rm occ}\overline n_k=E\), yields
\begin{equation}
 \log\frac{1+\overline n_k}{\overline n_k}
 =\alpha+\beta E_k^{\rm occ},
 \qquad
 \overline n_k
 =\frac{1}{e^{\beta(E_k^{\rm occ}-\mu)}-1},
 \qquad
 \mu=-\frac{\alpha}{\beta}.
 \label{eq:variacion-be-rev3}
\end{equation}
This is the Bose--Einstein occupation of
\eqref{eq:factorizacion-be-rev3}. The condition
\(\mu<\inf_kE_k^{\rm occ}\) ensures positivity of the denominator and
convergence of each modal factor.

In the finite instance constructed, the two constraints determine
\[
 (\beta_{\rm FD},\mu_{\rm FD})
 =(0.15307011353415928,\,4.502865197597043),
\]
\[
 (\beta_{\rm BE},\mu_{\rm BE})
 =(0.05456727960068316,\,-15.92569781906737).
\]
If \(m_k^{\rm mic}\) denotes the mean occupancy obtained by exact
microcanonical enumeration and \(m_k^{\rm eq}\) the evaluation of the corresponding
distribution, the errors
\(\sum_k|m_k^{\rm mic}-m_k^{\rm eq}|\) are
\[
 0.09897380039002812\quad\hbox{and}\quad
 0.9190625715582017,
\]
respectively. The enumeration traverses all occupation quadruples,
computes their integer multiplicities, and solves the two constraint
equations. The two laws are realized on the same screen. The physical choice
between them requires a grading and a Hamiltonian compatible with the
trajectory representation. The thermal parameter will subsequently be typed
as \(\beta=(k_B\Theta)^{-1}\).

\section{Quantified convergence to the Maxwell--Boltzmann distribution}

The relation between the two quantum statistics and the classical regime can
be formulated without resorting to an informal approximation. For each mode, let
\[
 u_k:=e^{-\beta(E_k^{\rm occ}-\mu)}.
\]
The state occupations are written
\[
 \overline n_k^{\rm FD}=\frac{u_k}{1+u_k},
 \qquad
 \overline n_k^{\rm BE}=\frac{u_k}{1-u_k};
\]
the multiplicities \(g_k\) are incorporated subsequently through
\(\langle n_k\rangle=g_k\overline n_k\).

\begin{theorem}[Common classical limit of quantum statistics]
\label{thm:limite-clasico-fd-be-rev2}
If \(0\leq u_k\leq\delta<1\), then
\[
 \left|\overline n_k^{\rm FD}-u_k\right|\leq u_k^2,
 \qquad
 \left|\overline n_k^{\rm BE}-u_k\right|
 \leq\frac{u_k^2}{1-\delta}.
\]
Consequently, in the regime dilute \(u_k\to0\),
\[
 \boxed{
 \overline n_k^{\rm FD}\sim
 \overline n_k^{\rm BE}\sim
 e^{-\beta(E_k^{\rm occ}-\mu)}} ,
\]
which is the Maxwell--Boltzmann distribution per mode.
\end{theorem}

\begin{proof}
The exact identities
\[
 u_k-\frac{u_k}{1+u_k}=\frac{u_k^2}{1+u_k},
 \qquad
 \frac{u_k}{1-u_k}-u_k=\frac{u_k^2}{1-u_k}
\]
provide the bounds. For \(0<u_k\leq\delta\), dividing by \(u_k\) and
taking \(u_k\downarrow0\) yields
the asymptotic equivalence. The leading term of order \(u_k^2\) preserves the
physical difference: exclusion decreases fermionic occupation, whereas
accumulation increases bosonic occupation.
\end{proof}
